% Part: computability
% Chapter: recursive-functions
% Section: general-recursive-functions

\documentclass[../../../include/open-logic-section]{subfiles}

\begin{document}

\olfileid[fr]{cmp}{rec}{gen}
\olsection{Fonctions récursives générales}

Il existe une autre manière de définir une classe de fonctions
totales. Une fonction totale $f(x,\vec z)$ est dite
\emph{régulière} si, pour tout tuple d'entiers naturels
$\vec z$, il existe un~$x$ tel que $f(x,\vec z)=0$.
Les fonctions régulières sont donc précisément celles
auxquelles la recherche non bornée peut être appliquée
pour obtenir une fonction totale. On peut restreindre
la recherche non bornée à ces seules fonctions :

\begin{defn}
\ollabel{defn:general-recursive}
La classe des \emph{fonctions récursives générales}
est la plus petite classe de fonctions d'arités quelconques,
des entiers naturels dans les entiers naturels, qui contient
la fonction nulle, le successeur et les projections, et
qui est close par composition, récurrence primitive et
recherche non bornée appliquée aux fonctions
\emph{régulières}.
\end{defn}

Il est clair que toute fonction récursive générale est
totale. La différence entre la
\olref{defn:general-recursive} et la
\olref[par]{defn:recursive-fn} tient au fait que la seconde
définition autorise l'emploi de fonctions récursives
partielles aux étapes intermédiaires : seule la fonction
obtenue à la fin doit être totale. L'adjectif « générale »
est donc un héritage historique trompeur; en apparence,
la \olref{defn:general-recursive} est \emph{moins} générale
que la \olref[par]{defn:recursive-fn}. Heureusement, cette
différence n'est qu'apparente : les deux définitions
déterminent la même classe de fonctions.\footnote{Complément explicatif : si $f$ est récursive partielle et totale, sa forme normale fournit un indice $e$ tel que, pour tout $x$, il existe $s$ vérifiant $T(e,x,s)$. La fonction récursive primitive $1\tsub\Char{T}(e,x,s)$ est donc régulière lorsqu'on cherche en $s$. Sa minimisation, puis la composition avec $U$, définissent $f$ sans fonction partielle intermédiaire. L'inclusion inverse découle des opérations de clôture.}

\end{document}
